%!TEX root = ../main.tex

\begin{abstract}
{
\noindent
Microbial carbon utilisation is a foundational ecological phenotype, yet remains difficult to predict from genomes despite well-characterised pathways.
Machine-learning models often generalise poorly across datasets, a failure usually attributed to training-set size and taxonomic bias.
To test this, we integrated binary growth phenotypes for 819 strains across 240 carbon sources from four datasets.
Balanced accuracy fell from 0.75--0.90 within datasets to 0.52--0.84 across them, and restricting tests to close relatives recovered only 0.03, so aggregating mechanistically inconsistent genotype-phenotype relationships drove models towards dataset-correlated shortcuts.
Restricting training to concordant samples, where GapMind predictions matched experiment, doubled the carbon sources recovering known pathway genes across datasets, from six to twelve of fifteen, whereas matched random subsets did not.
Coherence mattered more than volume: adding over 8000 literature-curated genomes did not significantly exceed the concordant set.
These models were bounded specialists, rescuing mechanistic false negatives twice as often as false positives (44\% versus 19\%), mostly metabolic generalists.
A mechanism-free filter using phylogenetic agreement and experimental labels, where no mechanistic predictor exists, recovered part of the gain but not the recall advantage.
Model confidence defined a per-prediction applicability domain, and prioritising low-confidence genomes improved cross-dataset accuracy more than random or diversity-based sampling, especially for the weakest phenotype.
These findings recast generalisation in biological machine learning as a problem of label-mechanism agreement and applicability-domain definition, alongside data volume and algorithm choice, and suggest a workflow of concordant training, applicability detection, and confidence-guided experimentation.
}
\end{abstract}

\keywords{Microbial ecology, Comparative genomics, Machine learning, Genotype-phenotype mapping, Cross-dataset generalisation}
